% method.tex -- the proposed (finalized) method: KnowledgeLLM.
% Describes the model architecture only (no datasets / experimental protocol).
% Organization follows the two core contributions:
%   (1) knowledge-bank extraction, (2) multimodal knowledge reprogramming.

%-------------------------------------------------------------------------
\section{Method}
\label{sec:method}

% TODO: pipeline figure (figures/framework.pdf). Suggested layout: input segment
% -> [raw / PSD / sliding-FC sidecars] -> patch encoders -> reprogramming
% cross-attention -> knowledge bank + mapper -> frozen LLM stack -> head.
% \begin{figure*}[t]
%   \centering
%   \includegraphics[width=\textwidth]{figures/framework.pdf}
%   \caption{Overview of \modelname.}
%   \label{fig:framework}
% \end{figure*}

\subsection{Overview}
\label{sec:method-overview}

\modelname adapts a frozen large language model (LLM) to EEG-based
classification and regression with only a small set of trainable parameters.
It rests on two components. First, \emph{knowledge-bank extraction}: the
parametric knowledge of the frozen LLM is read out as compact, aspect-specific
sets of vocabulary tokens that characterize the EEG views the model consumes.
Second, \emph{multimodal knowledge reprogramming}: three complementary EEG
views are encoded into the LLM's token space by cross-attention over learnable
fusion tokens derived from the bank
\TODO{\cite{time-llm}}, so the frozen LLM can consume EEG tokens with its
standard inference procedure, without fine-tuning any of its weights.

As illustrated in Fig.~\ref{fig:framework} (TODO: add figure), the pipeline
consists of five stages: multi-view input representation, patch encoding,
knowledge-bank extraction, knowledge reprogramming, and a frozen LLM backbone
with a lightweight prediction head. Only the patch encoders, the knowledge
mappers, the reprogramming layers, and the final head are trained; the LLM
backbone remains frozen throughout.

\subsection{Knowledge-Bank Extraction}
\label{sec:method-bank}

The knowledge source of \modelname is the frozen LLM itself. For each EEG
view---equivalently, each \emph{aspect} $a \in \{\mathrm{fc}, \mathrm{psd},
\mathrm{raw}\}$---we prepare a short statement $\mathcal{T}_a$ describing
clinically established EEG knowledge of that aspect, and pass it through the
\emph{complete} frozen LLM (before layer truncation) in plain text form. From
the logits of the final layer, we take the last $K$ non-special positions;
per position, special tokens are masked to $-\infty$ and the top-$M$ tokens
are retained and re-normalized with temperature $\tau_k$. Tokens are then
aggregated across positions by their probability mass and sorted, yielding
the aspect bank

\begin{equation}
  \mathcal{B}_a = \big\{ (v_i, w_i) \big\}_{i=1}^{|\mathcal{B}_a|},
\end{equation}
i.e., a set of vocabulary token ids $v_i$ with normalized weights $w_i$.
Intuitively, this activation-based extraction reads out which vocabulary
tokens the LLM most strongly associates with each EEG aspect, condensing the
LLM's parametric knowledge into a compact, reusable token set. The bank
embeddings
$\mathbf{E}_a = [\mathbf{e}_{v_1}; \dots; \mathbf{e}_{v_{|\mathcal{B}_a|}}]
\in \mathbb{R}^{|\mathcal{B}_a| \times D}$
are looked up from the frozen input embedding matrix of the LLM ($D$ is the
LLM hidden size), so the bank is a projection of the LLM's knowledge onto its
own token space. The extraction runs offline once and is reused throughout
training; only afterwards are the first $N$ layers of the LLM retained for
the backbone, keeping the trainable pipeline light.

\subsection{Multimodal Knowledge Encoding}
\label{sec:method-encoding}

The second component couples the bank to the EEG signal through three
complementary views, each tokenized by a patch encoder and then
\emph{reprogrammed} against the bank into the LLM's hidden space.

\paragraph{Multi-view input representation.}
Given one EEG segment $\mathbf{x} \in \mathbb{R}^{C \times T}$ with $C$
channels and $T$ time samples, each channel is first z-scored over time to
stabilize the input scale before convolution. We then build three views
covering the temporal, spectral, and connectivity aspects of the signal: the
normalized \emph{raw waveform} itself; a \emph{relative power spectral
density (PSD)} per channel, estimated Welch-like with a Hann window and
overlapping segments, band-limited and normalized to unit total power,
$\mathbf{S} \in \mathbb{R}^{C \times F}$ with $F$ frequency bins; and
\emph{sliding-window dynamic functional connectivity (FC)}, where the segment
is split into $W$ sliding windows, the Pearson correlation matrix
$\mathbf{R}^{(w)} \in \mathbb{R}^{C \times C}$ is computed per window, and
only the upper-triangular part is kept, giving an edge vector
$\mathbf{e}^{(w)} \in \mathbb{R}^{E}$ with $E = C(C-1)/2$. Each view carries
its own natural ``sequence'' axis: time, frequency, and windows.

\paragraph{Patch encoding.}
Each view is tokenized by a lightweight patch encoder with shared design
\TODO{\cite{ofa}}. For a view with sequence length $T_m$ and $D_m$ input
rows, the encoder unfolds the sequence into patches of size $p_m$ with stride
$s_m$; replication padding is applied so that a final partial patch is
preserved. The flattened patches are mapped to $d$-dimensional tokens by a
convolutional value embedding with kernel size 3 and circular padding,
followed by LayerNorm and dropout:

\begin{equation}
  \mathbf{H}_m = g_m(\mathbf{v}_m) \in \mathbb{R}^{L_m \times d},
  \qquad L_m = \left\lfloor \frac{T_m + s_m - p_m}{s_m} \right\rfloor + 1 ,
\end{equation}
where $\mathbf{v}_m$ denotes the view input. We deliberately add
\emph{no} positional embedding at this stage: the frozen LLM applies its own
rotary positional encoding downstream, and a fixed positional term here
dominates the convolutional output at initialization, which we found leads
to degenerate (collapsed) representations. The three encoders share the
design but use view-specific patch geometry and are not weight-shared,
producing feature tokens $\mathbf{H}_{\mathrm{raw}}$,
$\mathbf{H}_{\mathrm{psd}}$, and $\mathbf{H}_{\mathrm{fc}}$ of lengths
$L_{\mathrm{raw}}$, $L_{\mathrm{psd}}$, and $L_{\mathrm{fc}} = W$,
respectively.

\paragraph{Knowledge reprogramming.}
For each aspect, the bank is compressed into $P$ learnable \emph{fusion
tokens} by a bias-free linear mapper over the bank-token axis:

\begin{equation}
  \mathbf{f}_a = \mathbf{W}_a \mathbf{E}_a \in \mathbb{R}^{P \times D},
  \qquad \mathbf{W}_a \in \mathbb{R}^{P \times |\mathcal{B}_a|}.
\end{equation}
The mapper is initialized by \emph{segmented averaging}: the $i$-th output
row averages the $i$-th contiguous group of bank tokens with uniform
weights, so the $P$ fusion tokens are distinct from the very start. We found
this essential---a collinear initialization lets all fusion tokens collapse
to the same vector early in training. Dropout is applied to the fusion
tokens on the training path.

The view features are then \emph{reprogrammed} into the LLM space by a
multi-head cross-attention layer per aspect, where the feature tokens act as
queries and the fusion tokens act as both keys and values. With $H$ heads and
head dimension $d_h = d / H$,

\begin{equation}
  \mathbf{q} = \mathbf{Q}_a \mathbf{H}_a, \qquad
  \hat{\mathbf{k}} = \mathrm{normalize}\big(\mathbf{K}_a \mathbf{f}_a\big),
\end{equation}
\begin{equation}
  \alpha_{l,p}^{(h)} = \mathrm{softmax}_p\!\left(
    \frac{\langle \mathbf{q}_{l,h}, \hat{\mathbf{k}}_{p,h} \rangle}
         {\sqrt{d_h}} \right),
\end{equation}
\begin{equation}
  \mathbf{z}_a = \frac{1}{H} \sum_{h=1}^{H} \sum_{p=1}^{P}
    \alpha_{l,p}^{(h)} \mathbf{f}_{a,p} \;\in \mathbb{R}^{L_a \times D}.
\end{equation}
The key side is unit-normalized (\emph{scaled-cosine attention}): the frozen
word-embedding scale is far smaller than the query scale, and without
normalization the inter-key differences shrink to $\sim\!10^{-4}$ in score
space, degenerating the attention into a uniform average. Each output token
$\mathbf{z}_a$ is thus a query-specific mixture of the $P$ knowledge fusion
tokens, expressed directly in the LLM's hidden space. The three aspects are
concatenated into the final token sequence
$\mathbf{z} \in \mathbb{R}^{L \times D}$ with total length
$L = \sum_a L_a$:

\begin{equation}
  \mathbf{z} = [\,\mathbf{z}_{\mathrm{fc}};\; \mathbf{z}_{\mathrm{psd}};\;
                 \mathbf{z}_{\mathrm{raw}}\,].
\end{equation}

\subsection{Frozen LLM Backbone and Training Objectives}
\label{sec:method-objectives}

\paragraph{Frozen LLM backbone and prediction head.}
The token sequence $\mathbf{z}$ is consumed by the first $N$ layers of the
frozen LLM backbone (in our instantiation, an 8B-parameter Llama-style model
\TODO{\cite{deepseek-r1}}), using rotary positions $0, \dots, L-1$ and the
LLM's standard causal attention. All backbone weights are frozen; only the
RMSNorm parameters of the retained layers remain trainable, following
\TODO{\cite{ofa}}, so that the norms can absorb the scale of the injected EEG
tokens and match the LLM's internal statistics:

\begin{equation}
  \mathbf{h} = \mathrm{LLM}_{1:N}(\mathbf{z}) \in \mathbb{R}^{L \times D}.
\end{equation}
The representation is mean-pooled over tokens and fed into a small
prediction head:

\begin{equation}
  \hat{\mathbf{y}} = \mathrm{Head}\Big(
    \mathrm{GELU}\big(\mathrm{LayerNorm}(\bar{\mathbf{h}})\big)\Big),
  \qquad \bar{\mathbf{h}} = \frac{1}{L} \sum_{l=1}^{L} \mathbf{h}_l ,
\end{equation}
with dropout before the final linear projection.

\paragraph{Training objectives.}
For classification, the primary objective is the cross-entropy loss
$\mathcal{L}_{\mathrm{CE}}(\hat{\mathbf{y}}, y)$. We add a supervised
contrastive regularizer \TODO{\cite{supcon}} that aligns the multi-view
representations of samples sharing a label: for each aspect, the encoder
tokens are mean-pooled and projected,
$\mathbf{u}_a = \mathrm{proj}(\frac{1}{L_a}\sum_l \mathbf{H}_{a,l})$, and an
InfoNCE loss is computed over the batch with same-label pairs as positives
(self-pairs excluded), with temperature $\tau_c$:

\begin{equation}
  \mathcal{L} = \mathcal{L}_{\mathrm{CE}} + \lambda \,\mathcal{L}_{\mathrm{contrast}},
\end{equation}
\begin{equation}
  \mathcal{L}_{\mathrm{contrast}} =
    - \frac{1}{|\mathcal{V}|} \sum_{i \in \mathcal{V}}
      \log \frac{\sum_{p \in \mathcal{P}(i)} e^{s_{i,p}/\tau_c}}
                 {\sum_{j \neq i} e^{s_{i,j}/\tau_c}},
\end{equation}
where $s_{i,j}$ is the similarity between the projected view features of
samples $i$ and $j$, $\mathcal{P}(i)$ the positive set of $i$, and
$\mathcal{V}$ the set of samples with a valid positive pair. This term
encourages the three views to agree on the task semantics before the LLM
consumes them. For regression tasks, the head is trained with the Huber loss
($\delta = 1$).
